\section{Soundness, faithfulness and the ceiling}\label{sec:ceiling}

Let $G$ be the gate of the formal model, run with no customer attached.

\begin{theorem}[Soundness reduces to the write]\label{thm:sound}
If $W$ is faithful and complete on $h$, then $G(W(\view(h)), o) = \exec$ implies
$\Perm(h, o) = 1$.
\end{theorem}
\begin{proof}
Let $G$ execute $o = (a, \theta)$ and let $C$ be the $\Au$ records requesting $a$. By faithfulness
each request in $C$ is an own-voice request of the customer. A request is still listed only if no
execution spent it, and by completeness every outstanding request is listed. So the requests of
$C$ are exactly $\Out_a(h)$, and (i) holds.

(ii) $G$ takes $o^*$ as the single object the requests of $C$ name, and $\Perm$ does the same.
If none is named, $G$ takes the single value under $\kappa_a(s_a)$ on $\Au$ records. By
faithfulness and completeness, these values are exactly $\Vals_{\kappa_a(s_a)}(h)$, so that value
is $\st_{\kappa_a(s_a)}(h)$.

(iii) Now take $p \neq s_a$ and $k = \kappa_a(p)$. If the requests in $\Lambda \subseteq C$ name
exactly one value, $G$ binds it, and $\Perm$ agrees, because every request of $C$ names $o^*$ or
nothing and so lies in $\Lambda$. If they name none, $G$ binds $v$ only when the values on record
under $k$, over all labels, are $\{v\}$ and some $\Au$ record carries $v$. Faithfulness gives
$v \in \Vals_k(h)$. Completeness gives $\Vals_k(h) \subseteq \{v\}$. Hence $\st_k(h) = v$.
\end{proof}

\begin{theorem}[Faithfulness is also necessary]\label{thm:necessary}
Let $W$ be local, and let it be faithful and complete on own-voice customer messages. Suppose $W$
writes, for some message $m$, an $\Au$ record $r$ that is not backed in the sense of
Definition~\ref{def:faithful}: a slot $(k, v)$ where $m$ asserts no value under $k$ in own
voice, or a request that $m$ does not make. Then there is a history of two messages on which $G \circ W$ executes
an operation with $\Perm = 0$.
\end{theorem}
\begin{proof}
\emph{The slot case.} Take an action $a$ with a parameter $p \neq s_a$ read from $k$. Let
$h' = (m, u)$ with $B = \emptyset$, where $u$ is an own-voice request for $a$ that names the object
and every parameter but $p$ and asserts nothing under $k$. By locality, $W$ writes $r$ in $h'$
exactly as before; it writes $u$ faithfully. The only value on record under $k$ is $v$, carried by
the $\Au$ record $r$, so $G$ binds $p = v$ and executes. But $\Vals_k(h') = \emptyset$, so
$\Perm = 0$.

\emph{The request case.} Let $u$ assert, in own voice and without requesting anything, the values
that $r$'s request names and one value for every other parameter of the requested action. $G$ is licensed by $r$ and binds
every argument, while $\Out_a(h') = \emptyset$.
\end{proof}

Theorems~\ref{thm:sound} and~\ref{thm:necessary} make the conditional in ``zero attacks under
correct labels'' exact. Under completeness, $G$ is sound on every continuation if and only if the
write is faithful. Taking expectations over histories gives a bound that can be measured one part
at a time:
\[
\mathrm{ASR}(G \circ W) \;\leq\; \Pr\big[\,W \text{ not faithful or not complete on } H^{-}\big].
\]

\begin{theorem}[The ceiling]\label{thm:ceiling}
(a) The write rule of the formal model satisfies $\ell(r) \preceq \ceil{\chi(r)}$ for every record,
and the store refuses any other record.
(b) Let $W$ be local. If $W$ writes an $\Au$ record carrying a slot or a request for some message
$(\chi, x)$ with $\chi \neq \cU$, then under Assumption~\ref{as:adv}(a) $W$ is unfaithful on some
history, and $G \circ W$ executes an unpermitted operation on some history. The same holds for
$\chi = \cU$ under Assumption~\ref{as:adv}(b), on the history in which $x$ is a quotation.
\end{theorem}
\begin{proof}
(a) Read off the cases of the write rule. (b) By Assumption~\ref{as:adv}(a) the message
$(\chi, x)$ occurs in some history. By locality, the same $\Au$ record is written there. A message
on $\chi \neq \cU$ is never in $\Own(h)$, so the record is unbacked, and
Theorem~\ref{thm:necessary} applies. For $\cU$, Assumption~\ref{as:adv}(b) places $x$ in a history
with $\omega = \quoted$, where the record is again unbacked.
\end{proof}

So on every channel but the customer's, the ceiling is the highest label a local write can safely
give. What a message says can lower its label but never raise it. On the customer's channel even
$\Au$ cannot be certified from the text. There it rests on a model's reading of the voice, which is
the subject of Section~\ref{sec:speech}.

\begin{corollary}[Grants]\label{cor:grant}
A grant arriving on $\cT$ cannot license an action, however it is verified. $\Au$ on such a
record is unbacked on the history in which the customer never asked. A standing mandate can
license only if it arrives on a channel whose messages are the customer's own voice by
construction, one that carries the customer's verifiable signature. That is a new channel with
ceiling $\Au$, not a property of the bank system that relays it.
\end{corollary}

\begin{proposition}[Lowering is safe only away from the customer]\label{prop:lower}
Lowering the label of a record not backed by an own-voice message preserves faithfulness and
completeness wherever they held, and so preserves soundness. Lowering the label of a backed record breaks completeness, and
$G \circ W$ can then execute an unpermitted operation.
\end{proposition}
\begin{proof}
The first claim holds because such a record is neither needed for completeness nor allowed to be
$\Au$ by faithfulness. For the second, consider
$h = \big(\text{``Send 500 EUR to Anna''},\ \text{``No wait, make it 300''}\big)$, both own-voice
requests for a transfer, the first naming $500$ and the second $300$. The named amounts are
$\{500, 300\}$, so every transfer has $\Perm = 0$. Suppose $Q$ wrongly calls the second message a
quotation. It is then stored as $\Un$ with no requests; the only request left names $500$, and $G$
sends 500 EUR. The same happens to the object when a lowered customer message carries a second
object: step~1 of $G$ reads objects from $\Au$ records only.
\end{proof}

\begin{remark}
So a false ``quote'' on a customer message is not harmless. It can hide a correction.
Proposition~\ref{prop:lower} makes precise the rule that a label may be lowered but never
raised: lowering is free for everything that is not the customer's own words. For
arguments, step~2 already reads values over all labels, and a lowered value still blocks. Step~1
and the request list do not. Whether to close this is a design decision; it would cost utility
only.
\end{remark}
